% GENERATED FILE. Edit canonical lessons, then run npm run generate:books.
% canonical-source-hash: fnv1a64:0594464b5fe9b290
% canonical-lessons: GE-C381-vorschrift, GE-C381-vorschreiben, GE-C381-beachten, GE-C381-einhalten, GE-C381-verpflichtet

\chapter{Vorschriften}
\label{ch:ge-vorschriften}
\hlchaptermodality{381}

\begin{chapteropening}
\textbf{By the end of this chapter:} \emph{I can talk about regulations: what is laid down, what we must observe and keep to, and what we are obliged to do.}

Complete the last lesson of chapter 381: I can talk about regulations: what is laid down, what we must observe and keep to, and what we are obliged to do.
\end{chapteropening}

\section[die Vorschrift]{die Vorschrift — the regulation}

\label{lesson:GE-C381-vorschrift}



\begin{quote}
Before the new one: say the German for the registration, then the German for to withdraw.
\end{quote}

\subsection*{You'll want to know: die Vorschrift}
\textbf{die Vorschrift} — \textquotedblleft{}the regulation\textquotedblright{}.

\begin{quote}
Nach der Vorschrift muss man im Labor eine Schutzbrille tragen.
\end{quote}

\begin{quote}
\emph{According to the regulations you have to wear goggles in the lab.}
\end{quote}

Plural: \emph{die Vorschriften}. \emph{Dienst nach Vorschrift} — work-to-rule.

\begin{cousinweb}
\emph{Vorschrift} goes with \emph{vorschreiben} \textquotedblleft{}to lay down\textquotedblright{}, built on \emph{Schrift} \textquotedblleft{}writing\textquotedblright{}, from Old High German \emph{scrift} and Latin \emph{scrīptum} (EtymDB): a rule written out in front of you.
\end{cousinweb}

\begin{grammarlens}[title={What you've built}]
One more word for telling, explaining and describing at B1.
\end{grammarlens}

\subsection*{Your turn}
\begin{itemize}
\raggedright
  \item \emph{Say it:} \emph{die Vorschrift}
  \item \emph{Say it:} \emph{die Vorschrift}, once more
  \item \emph{Say it:} a sentence of your own with \emph{die Vorschrift}
  \item \emph{Recall:} say the German for the registration, then the German for to withdraw, then say \emph{die Vorschrift} again
\end{itemize}

\begin{tcolorbox}[breakable,colback=teal!4,colframe=teal!35!black,title={Before you move on}]
What does die Vorschrift mean? (\textquotedblleft{}The regulation\textquotedblright{}.) Say it once more.
\end{tcolorbox}
\section[vorschreiben]{vorschreiben — to lay down (a rule), to prescribe}

\label{lesson:GE-C381-vorschreiben}



\begin{quote}
Before the new one: say the German for to withdraw, then the German for the regulation.
\end{quote}

\subsection*{You'll want to know: vorschreiben}
\textbf{vorschreiben} — \textquotedblleft{}to lay down (a rule), to prescribe\textquotedblright{}.

\begin{quote}
Das Gesetz schreibt vor, dass Kinder im Auto einen Kindersitz brauchen.
\end{quote}

\begin{quote}
\emph{The law lays down that children need a child seat in the car.}
\end{quote}

Separable and strong: \emph{es schreibt vor}, \emph{es hat vorgeschrieben}. \emph{Du kannst mir nichts vorschreiben} — you cannot tell me what to do.

\begin{cousinweb}
\emph{vorschreiben} is built from \emph{vor} \textquotedblleft{}before\textquotedblright{} and \emph{schreiben}, from Old High German \emph{scrīban}, a very early borrowing from Latin \emph{scrībere} (EtymDB); EtymDB lists English \emph{scribe} as a cousin.
\end{cousinweb}

\begin{grammarlens}[title={What you've built}]
One more word for telling, explaining and describing at B1.
\end{grammarlens}

\subsection*{Your turn}
\begin{itemize}
\raggedright
  \item \emph{Say it:} \emph{vorschreiben}
  \item \emph{Say it:} \emph{vorschreiben}, once more
  \item \emph{Say it:} a sentence of your own with \emph{vorschreiben}
  \item \emph{Recall:} say the German for to withdraw, then the German for the regulation, then say \emph{vorschreiben} again
\end{itemize}

\begin{tcolorbox}[breakable,colback=teal!4,colframe=teal!35!black,title={Before you move on}]
What does vorschreiben mean? (\textquotedblleft{}To lay down (a rule), to prescribe\textquotedblright{}.) Say it once more.
\end{tcolorbox}
\section[beachten]{beachten — to observe, to pay attention to}

\label{lesson:GE-C381-beachten}



\begin{quote}
Before the new one: say the German for the regulation, then the German for to lay down.
\end{quote}

\subsection*{You'll want to know: beachten}
\textbf{beachten} — \textquotedblleft{}to observe, to pay attention to\textquotedblright{}.

\begin{quote}
Viele Autofahrer beachten das Tempolimit nicht.
\end{quote}

\begin{quote}
\emph{Many drivers do not observe the speed limit.}
\end{quote}

Inseparable: \emph{er beachtet, er hat beachtet}. \emph{Etwas nicht beachten} — to ignore something.

\begin{cousinweb}
\emph{beachten} is inherited from Middle High German \emph{beahten} and Old High German \emph{biahtōn} (EtymDB), built on \emph{Acht} \textquotedblleft{}attention\textquotedblright{}, Old High German \emph{ahta} (EtymDB): you pay attention.
\end{cousinweb}

\begin{grammarlens}[title={What you've built}]
One more word for telling, explaining and describing at B1.
\end{grammarlens}

\subsection*{Your turn}
\begin{itemize}
\raggedright
  \item \emph{Say it:} \emph{beachten}
  \item \emph{Say it:} \emph{beachten}, once more
  \item \emph{Say it:} a sentence of your own with \emph{beachten}
  \item \emph{Recall:} say the German for the regulation, then the German for to lay down, then say \emph{beachten} again
\end{itemize}

\begin{tcolorbox}[breakable,colback=teal!4,colframe=teal!35!black,title={Before you move on}]
What does beachten mean? (\textquotedblleft{}To observe, to pay attention to\textquotedblright{}.) Say it once more.
\end{tcolorbox}
\section[einhalten]{einhalten — to keep to, to comply with}

\label{lesson:GE-C381-einhalten}



\begin{quote}
Before the new one: say the German for to lay down, then the German for to observe.
\end{quote}

\subsection*{You'll want to know: einhalten}
\textbf{einhalten} — \textquotedblleft{}to keep to, to comply with\textquotedblright{}.

\begin{quote}
Wir müssen die Frist einhalten.
\end{quote}

\begin{quote}
\emph{We have to meet the deadline.}
\end{quote}

Separable and strong: \emph{er hält ein}, \emph{er hat eingehalten}. \emph{Ein Versprechen einhalten} — to keep a promise.

\begin{cousinweb}
\emph{einhalten} is built from \emph{ein} and \emph{halten} \textquotedblleft{}to hold\textquotedblright{}, inherited from Proto-Germanic \emph{haldaną} (EtymDB): you hold to what was agreed.
\end{cousinweb}

\begin{grammarlens}[title={What you've built}]
One more word for telling, explaining and describing at B1.
\end{grammarlens}

\subsection*{Your turn}
\begin{itemize}
\raggedright
  \item \emph{Say it:} \emph{einhalten}
  \item \emph{Say it:} \emph{einhalten}, once more
  \item \emph{Say it:} a sentence of your own with \emph{einhalten}
  \item \emph{Recall:} say the German for to lay down, then the German for to observe, then say \emph{einhalten} again
\end{itemize}

\begin{tcolorbox}[breakable,colback=teal!4,colframe=teal!35!black,title={Before you move on}]
What does einhalten mean? (\textquotedblleft{}To keep to, to comply with\textquotedblright{}.) Say it once more.
\end{tcolorbox}
\section[verpflichtet]{verpflichtet — obliged, required}

\label{lesson:GE-C381-verpflichtet}



\begin{quote}
Before the new one: say the German for to observe, then the German for to keep to.
\end{quote}

\subsection*{You'll want to know: verpflichtet}
\textbf{verpflichtet} — \textquotedblleft{}obliged, required\textquotedblright{}.

\begin{quote}
Der Arbeitgeber ist verpflichtet, den Lohn pünktlich zu zahlen.
\end{quote}

\begin{quote}
\emph{The employer is obliged to pay wages on time.}
\end{quote}

\emph{Zu etwas verpflichtet sein} — to be obliged to do something. From \emph{verpflichten} — to oblige.

\begin{cousinweb}
\emph{verpflichtet} is the participle of \emph{verpflichten}, built on \emph{Pflicht} \textquotedblleft{}duty\textquotedblright{}, inherited from Old High German \emph{pfliht} and Proto-Germanic \emph{plihtiz} (EtymDB).
\end{cousinweb}

\begin{grammarlens}[title={What you've built}]
One more word for telling, explaining and describing at B1.
\end{grammarlens}

\subsection*{Your turn}
\begin{itemize}
\raggedright
  \item \emph{Say it:} \emph{verpflichtet}
  \item \emph{Say it:} \emph{verpflichtet}, once more
  \item \emph{Say it:} a sentence of your own with \emph{verpflichtet}
  \item \emph{Recall:} say the German for to observe, then the German for to keep to, then say \emph{verpflichtet} again
\end{itemize}

\begin{tcolorbox}[breakable,colback=teal!4,colframe=teal!35!black,title={Before you move on}]
What does verpflichtet mean? (\textquotedblleft{}Obliged, required\textquotedblright{}.) Say it once more.
\end{tcolorbox}
